% ============================================================
% Capítulo: Declaraciones (Bloque B.15)
% ============================================================
\chapter{Declaraciones}
\label{cap:declaraciones}

Este capítulo reúne las declaraciones formales que la guía exige antes de
las referencias bibliográficas. Sigue el mismo criterio de honestidad
metodológica que el resto del documento: los datos verificables en el
propio repositorio se presentan como tales, y los que exigen una
decisión que solo puede tomar el equipo (financiamiento, conflictos de
interés, distribución exacta de contribuciones) se dejan marcados de
forma explícita como pendientes, en vez de decidirse unilateralmente en
este documento.

\section{Disponibilidad de código}
\label{sec:decl-codigo}

El código fuente completo de SGB-SaaS (backend Spring Boot, frontend
Angular, scripts de medición, infraestructura Docker) está disponible
públicamente bajo licencia MIT (\texttt{LICENSE}, raíz del repositorio)
en:

\begin{center}
\url{https://github.com/mloorm14/sgb-saas}
\end{center}

El software citable de esta entrega corresponde al tag \texttt{v1.0.0}
(ver portada de este documento para el hash de commit exacto, fijado al
cierre). Los metadatos de citación siguen el formato Citation File
Format 1.2.0 (\texttt{CITATION.cff}, raíz del repositorio), con ORCID
verificado de los 3 autores.

\section{Disponibilidad de datos}
\label{sec:decl-datos}

Los datos empíricos crudos que sustentan el \autoref{cap:resultados}
(salidas de k6, reportes JaCoCo en XML/CSV/HTML, resultados de
Lighthouse, evidencia de auditoría OWASP control por control) están
versionados junto con el código fuente en \texttt{docs/mediciones/}, en
el mismo repositorio citado en la
\autoref{sec:decl-codigo}. \texttt{LICENSE} cubre el repositorio
completo (\emph{"the Software"}, sin excluir explícitamente los datos de
\texttt{docs/}), por lo que estos datos quedan bajo la misma licencia
MIT que el software -- el equipo no ha declarado hasta la fecha una
licencia de datos separada (p.~ej.\ CC-BY) para este material, algo que
podría formalizarse en una futura revisión de \texttt{LICENSE} si el
equipo lo considera necesario.

A diferencia del software (con DOI de Zenodo, ver más abajo), este
conjunto de datos empíricos \textbf{no cuenta todavía con un DOI propio
independiente} -- no se ha creado un archivado separado en un
repositorio de datos (Zenodo permite archivar datasets con su propio
DOI, distinto del DOI del software). El campo queda como
\emph{placeholder} explícito hasta que el equipo decida si amerita un
archivado independiente o si basta con el DOI único del software
(criterio también pendiente de decisión del equipo):

\begin{center}
\texttt{<DOI-DATASET-V1.0.0-SI-APLICA>}
\end{center}

Los formularios de consentimiento informado firmados por participantes
reales de las pruebas de usabilidad (cuando se ejecuten,
\autoref{sec:tf-sus}) están explícitamente excluidos de este
repositorio y de cualquier archivado público, por la razón declarada en
la \autoref{sec:decl-etica}.

\section{Disponibilidad de materiales}
\label{sec:decl-materiales}

Los instrumentos y materiales usados para producir la evidencia de este
documento están versionados como parte del propio repositorio de
código, no como capturas de pantalla ni archivos externos no
trazables:

\begin{itemize}[leftmargin=1.2cm]
  \item \textbf{Colección de peticiones HTTP}: \texttt{docs/postman/coleccion.json},
    con 66 peticiones reales que cubren los módulos de
    Auth, Libros, Préstamos, Reservaciones, Multas, Configuración,
    Credencial QR, Notificaciones, Usuarios (Admin), Auditoría y
    Chatbot, incluyendo casos de error (401/403/404/422/423/429). Nota
    de honestidad: \texttt{docs/observaciones/OBSERVACIONES.md} (OBS-03)
    registra 39 peticiones como la cifra vigente en el commit
    \texttt{c6b372c} (cierre de la Tercera Entrega); el conteo real
    verificado directamente sobre el archivo actual para este capítulo
    es 66 -- la colección creció junto con los 8 módulos de dominio
    mergeados durante esta entrega, y esa cifra antigua no se actualizó
    en su momento en \texttt{OBSERVACIONES.md}.
  \item \textbf{Scripts de carga (k6)}: \texttt{k6/}, con semilla fija
    documentada (\autoref{sec:mm-semillas}) para reproducibilidad.
  \item \textbf{Script de análisis estadístico}: \texttt{scripts/perf-analysis.py}
    (bootstrap de IC~95\,\% sobre p95, semilla fija \texttt{BOOTSTRAP\_SEED = 42}).
  \item \textbf{Configuración de Lighthouse CI}: \texttt{frontend-angular/lighthouserc.js},
    con el perfil móvil/Slow~4G ya parametrizado y comentado, listo para
    ejecutarse contra el stack Docker real.
  \item \textbf{Protocolo de usabilidad SUS}: plantilla de consentimiento
    informado en \texttt{docs/etica/consentimientos/plantilla.md}
    (\autoref{sec:tf-sus}).
\end{itemize}

\section{Declaración ética}
\label{sec:decl-etica}

El tratamiento de datos personales y el protocolo de consentimiento
informado para las pruebas de usabilidad SUS están documentados en
detalle en \texttt{docs/etica/ETHICS.md}, verificado directamente para
este capítulo. En síntesis:

\begin{itemize}[leftmargin=1.2cm]
  \item Los datos de ejemplo del sistema (\texttt{db/seed.sql}) usan
    metadatos bibliográficos públicos de libros reales publicados (ISBN,
    título, autor, editorial); ninguna contraseña se almacena en texto
    plano (hash BCrypt, costo 12); se verificó por auditoría de código
    que ningún endpoint expone el campo \texttt{passwordHash}.
  \item Los participantes de las pruebas SUS (cuando se ejecuten,
    \autoref{sec:tf-sus}) firman un consentimiento informado bajo la
    plantilla ya versionada, que declara explícitamente el propósito de
    la prueba, los datos recogidos (sin datos identificables, solo un
    código de participante \texttt{P01}, \texttt{P02}, ...), el
    mecanismo de anonimización, y el derecho a retirarse en cualquier
    momento sin consecuencias.
  \item Los formularios de consentimiento ya firmados por participantes
    reales \textbf{no se suben a este repositorio público bajo ninguna
    forma} -- se conservan, según el texto verificado de
    \texttt{ETHICS.md}, "por el equipo en un medio separado (físico o
    digital con acceso restringido)". Nota de precisión: este documento
    inicialmente iba a describir ese medio como "Drive institucional con
    acceso restringido al docente", pero \texttt{ETHICS.md} no especifica
    ese detalle -- solo declara que el medio es separado del repositorio
    y de acceso restringido, sin nombrar la herramienta concreta. Se
    reporta aquí lo que el archivo realmente dice, no el detalle asumido
    inicialmente.
\end{itemize}

\section{Contribuciones de autoría (taxonomía CRediT)}
\label{sec:decl-credit}

\textbf{Nota de honestidad metodológica.} La distribución de roles que
sigue es una \textbf{propuesta borrador}, construida a partir de
evidencia real verificable en el repositorio (autoría de commits por
área funcional, vía \texttt{git shortlog}, y asignaciones ya registradas
en \texttt{docs/observaciones/OBSERVACIONES.md}) -- \textbf{no es una
declaración definitiva}. Queda sujeta a revisión y confirmación
explícita por los 3 integrantes antes del cierre de la entrega; ningún
integrante debe asumir que esta tabla refleja un acuerdo ya alcanzado.

\begin{table}[h]
\centering
\small
\begin{tabular}{@{}p{4.6cm}ccc@{}}
\toprule
\textbf{Rol CRediT} & \textbf{Cajas} & \textbf{Loor} & \textbf{Panamá} \\
\midrule
Conceptualization        & \checkmark & \checkmark & \checkmark \\
Data curation             &            & \checkmark &            \\
Formal analysis           &            & \checkmark &            \\
Funding acquisition       & \multicolumn{3}{c}{No aplica -- sin financiamiento externo declarado} \\
Investigation              & \checkmark & \checkmark & \checkmark \\
Methodology                & \checkmark & \checkmark &            \\
Project administration     &            & \checkmark &            \\
Resources                  & \checkmark & \checkmark & \checkmark \\
Software                   & \checkmark & \checkmark & \checkmark \\
Supervision                & \multicolumn{3}{c}{Ver nota debajo de la tabla} \\
Validation                  & \checkmark & \checkmark &            \\
Visualization               &            & \checkmark &            \\
Writing -- original draft    &            & \checkmark &            \\
Writing -- review \& editing  & \checkmark & \checkmark & \checkmark \\
\bottomrule
\end{tabular}
\caption{Propuesta borrador de distribución de roles CRediT, pendiente de confirmación del equipo.}
\label{tab:decl-credit}
\end{table}

Justificación de la propuesta, con evidencia concreta verificada para
este capítulo:

\begin{itemize}[leftmargin=1.2cm]
  \item \textbf{Cajas}: autor del mayor volumen de commits del proyecto
    (194 de $\sim$390 commits con identidad atribuible, contando las
    variantes de firma \texttt{Theirvin1}/\texttt{TheIrvin}/\texttt{Irvin}),
    concentrados en el backend Spring Boot -- servicios, controladores,
    la migración forzada de \texttt{@Procedure} a \texttt{@Query} nativa
    por el defecto de Hibernate (\autoref{cap:discusion}), y la mayoría
    de los 203 tests automatizados del backend. Autor de los 8 módulos
    de dominio mergeados en esta entrega. Asignado también a la
    auditoría ZAP baseline pendiente (\autoref{sec:tf-zap}).
  \item \textbf{Loor}: autor de la mayoría de los commits de
    \texttt{docs/} (este informe, arquitectura C4, ADRs, bibliografía
    verificada contra Crossref), de los scripts de medición
    (\texttt{k6/}, \texttt{scripts/perf-analysis.py}) y del análisis
    estadístico del \autoref{cap:resultados}. 150 commits (variantes
    \texttt{Marlon Loor}/\texttt{mloorm14}/\texttt{Loor Marlon}).
  \item \textbf{Panamá}: autor de los componentes de frontend de
    Préstamos, Reservaciones y Multas (\texttt{frontend-angular/}) y de
    varias historias de usuario en \texttt{docs/requisitos/}. 44
    commits. Asignado a la ejecución de las pruebas SUS pendientes
    (OBS-08, \autoref{sec:tf-sus}).
\end{itemize}

\textbf{Nota sobre "Supervision".} La guía CRediT reserva este rol
típicamente para quien dirige el trabajo sin ser autor operativo directo
del código -- en este PFC, ese rol lo ejerce el docente-director (Dr.
Gleiston Cicerón Guerrero Ulloa, ver portada), no uno de los 3
integrantes-autores. Se deja fuera de la tabla de columnas por
integrante para no asignarlo incorrectamente a ninguno de los 3, y se
señala aquí como punto abierto: el equipo debe decidir si desea listar
al docente-director bajo este rol en la versión final de esta
declaración.

\section{Conflictos de interés}
\label{sec:decl-conflictos}

% PLACEHOLDER: el equipo declara si existe o no algún conflicto de
% interés (relación laboral, financiera o personal con alguna entidad
% que pudiera influir en los resultados reportados) -- completar antes
% de la entrega final. Dado que este es un Proyecto de Fin de Curso
% (PFC) académico sin patrocinio externo conocido, lo esperable es una
% declaración de ausencia de conflictos, pero esa afirmación debe
% hacerla el equipo explícitamente, no asumirse aquí por el redactor de
% este capítulo.

\texttt{[PLACEHOLDER: completar antes de la entrega final -- ver nota
de honestidad arriba. El equipo debe declarar explícitamente presencia o
ausencia de conflictos de interés.]}

\section{Financiamiento}
\label{sec:decl-financiamiento}

% PLACEHOLDER: el equipo declara si hubo o no financiamiento externo
% (beca, subvención, patrocinio de una entidad) para este proyecto --
% completar antes de la entrega final. No se encontró evidencia de
% financiamiento en el repositorio (sin archivo de patrocinio, sin
% mención en README/CITATION.cff), pero la ausencia de evidencia en el
% repositorio no es lo mismo que una declaración formal del equipo.

\texttt{[PLACEHOLDER: completar antes de la entrega final -- el equipo
declara si hubo o no financiamiento externo. No se encontró evidencia de
financiamiento en el repositorio al momento de escribir este capítulo,
pero la declaración formal corresponde al equipo, no a una inferencia
por ausencia de evidencia.]}

\section{Uso de inteligencia artificial generativa}
\label{sec:decl-ia}

Este proyecto usó de forma extensiva el asistente Claude (Anthropic) a
lo largo de todo el ciclo de desarrollo, con dos roles diferenciados:

\begin{itemize}[leftmargin=1.2cm]
  \item \textbf{Claude}, en conversación directa, como orquestador y
    revisor: definición de tareas técnicas concretas, redacción asistida
    de este mismo documento académico (incluyendo este capítulo),
    verificación cruzada de referencias bibliográficas contra la API de
    Crossref (\autoref{cap:trabajos-relacionados}), y análisis de datos
    empíricos ya generados (interpretación de salidas de k6, JaCoCo,
    Lighthouse y auditoría OWASP para el \autoref{cap:resultados}).
  \item \textbf{Claude Code}, como agente ejecutor de las tareas
    técnicas derivadas de esas conversaciones: generación y edición de
    código fuente (backend, frontend, scripts), ejecución de comandos
    (compilaciones, builds Docker, corridas de tests), y operaciones de
    control de versiones (commits, push).
\end{itemize}

En ambos casos, cada cambio propuesto por la IA -- ya fuera código,
texto de este informe, o un comando a ejecutar -- pasó por revisión
humana explícita del integrante que dirigía esa sesión de trabajo antes
de aplicarse: los commits del repositorio no se generaron ni se
subieron de forma autónoma sin intervención humana en el punto de
decisión. Este documento declara este uso con el mismo criterio de
honestidad metodológica aplicado al resto de sus afirmaciones: ni se
minimiza (la IA participó de forma sustancial, no marginal, en la
redacción de la documentación y en buena parte de la implementación) ni
se exagera (las decisiones de diseño, arquitectura, alcance y los
juicios sobre qué evidencia era suficiente o qué limitación declarar
correspondieron al equipo humano, no a la IA de forma autónoma).

\section{Agradecimientos}
\label{sec:decl-agradecimientos}

% PLACEHOLDER: opcional. El equipo puede mencionar aquí al
% docente-director u otras personas/instituciones que colaboraron
% durante el desarrollo de este proyecto, si así lo desea.

\texttt{[PLACEHOLDER: sección opcional -- el equipo puede mencionar aquí
al docente-director u otras personas/instituciones que colaboraron, si
lo desea.]}
